\section*{第 \ref{ch:g8:cosine} 章 --- 直角三角形与余弦}

\begin{solution}{exo:g8:cosine:1}
直角所对的中线等于斜边的一半，即 $6.5$ 厘米
（\cref{thm:g8:cosine:circle}）。外接圆的圆心位于斜边中点，
半径为 $6.5$ 厘米。
\end{solution}

\begin{solution}{exo:g8:cosine:2}
点 $A$ 处的角为直角（$[BC]$ 为直径，
由\cref{thm:g8:cosine:circle}的第 2 点）。测量值满足
$AB^2 + AC^2 \approx 64 = BC^2$，在测量精度范围内。
\end{solution}

\begin{solution}{exo:g8:cosine:3}
斜边为 $[EF]$（直角 $D$ 所对的边）。
与 $\widehat E$ 相邻的边为 $[ED]$，对边为 $[DF]$。于是
$\cos\widehat E = \dfrac{ED}{EF}$。
\end{solution}

\begin{solution}{exo:g8:cosine:4}
\emph{1.} 邻边 $= 10 \cos 35^\circ \approx 8.2$。

\emph{2.} 斜边 $= \dfrac{6}{\cos 50^\circ} \approx
\dfrac{6}{0.643} \approx 9.3$。
\end{solution}

\begin{solution}{exo:g8:cosine:5}
$\cos\theta = \frac{8}{10} = 0.8$，于是
$\theta = \cos^{-1}(0.8) \approx 37^\circ$。
\end{solution}

\begin{solution}{exo:g8:cosine:6}
$\cos\theta$ 等于直角边除以斜边，而斜边
是直角三角形中\emph{最长}的边，因此该商
总小于 $1$。若取值为 $1.2$，则直角边长于
斜边 --- 这是不可能的。
\end{solution}

\begin{solution}{exo:g8:cosine:7}
水平跨度 $= 25 \cos 12^\circ \approx 25 \times 0.978 \approx
24.5$ 米。
\end{solution}

\begin{solution}{exo:g8:cosine:8}
余弦随角增大而减小
（见\cref{prop:g8:cosine:bounds}）：
\[
\cos 70^\circ < \cos 45^\circ < \cos 20^\circ .
\]
（计算器：$0.342 < 0.707 < 0.940$。）
\end{solution}

\begin{solution}{exo:g8:cosine:9}
地面距离与 $\theta$ 相邻：
$\cos 68^\circ = \frac{1.2}{L}$，于是
$L = \frac{1.2}{0.375} = 3.2$ 米。高度（由毕达哥拉斯定理）：
$\sqrt{3.2^2 - 1.2^2} = \sqrt{10.24 - 1.44} = \sqrt{8.8} \approx 3.0$
米。
\end{solution}

\begin{solution}{exo:g8:cosine:10}
\emph{1.} 点 $A$ 在以 $[BC]$ 为直径的圆上，$A$ 处为直角，
由\cref{thm:g8:cosine:circle}可知。

\emph{2.} $BC = 9$ 厘米（半径的两倍）。由毕达哥拉斯定理：
$AC = \sqrt{9^2 - 5.4^2} = \sqrt{81 - 29.16} = \sqrt{51.84} = 7.2$
厘米。且 $\cos\widehat B = \dfrac{AB}{BC} = \dfrac{5.4}{9} = 0.6$。
\end{solution}

\begin{solution}{exo:g8:cosine:11}
边长为 $1$ 的等边三角形中，自一个顶点
落在对边的\emph{中点}上（该直线为
对称轴）。三角形的一半是直角三角形，斜边为 $1$，
（即一条完整的边），底边端点处的角为 $60^\circ$，
邻边为 $\frac12$（底边的一半）。因此
\[
\cos 60^\circ = \frac{1/2}{1} = \frac12 .
\]
\end{solution}

\begin{solution}{pb:g8:cosine:1}
\textbf{1.} 在三角形 $ABH$ 中，$H$ 为直角，
由\cref{thm:g7:triangles:sum}得
$\widehat B + 90^\circ + \widehat{BAH} = 180^\circ$，于是
$\widehat{BAH} = 90^\circ - \widehat B$。在三角形 $ABC$ 中，
$A$ 为直角：
$\widehat B + \widehat C + 90^\circ = 180^\circ$，于是
$\widehat C = 90^\circ - \widehat B$。因此
$\widehat{BAH} = \widehat{ACB}$：两个小三角形的
角 $\widehat B$、$\widehat C$、$90^\circ$ 与原三角形相同。

\textbf{2.} 在 $ABC$ 中（$A$ 为直角），与 $\widehat B$ 相邻的边为
$AB$，斜边为 $BC$：
$\cos\widehat B = \frac{AB}{BC}$。在 $ABH$ 中（$H$ 为直角），
与 $\widehat B$ 相邻的边为 $BH$，斜边为
$AB$：$\cos\widehat B = \frac{BH}{AB}$。该比值只依赖于
角（见\cref{def:g8:cosine:def}），于是
\[
\frac{AB}{BC} = \frac{BH}{AB} ;
\]
两边同乘 $BC \times AB$（
由\cref{thm:g8:equations:balance}）：
$AB^2 = BH \times BC$。

\textbf{3.} 三角形 $ABC$（$A$ 为直角）与 $ACH$
（$H$ 为直角）有公共角 $\widehat C$。两者中，邻边
与斜边之比分别为：
\[
\cos\widehat C = \frac{AC}{CB} = \frac{CH}{AC}
\qquad\Longrightarrow\qquad
AC^2 = CH \times CB .
\]

\textbf{4.} 代入 $AB^2 = BH \times BC$ 得 $9 = BH \times 5$，故
$BH = 1.8$；又 $16 = CH \times 5$ 得 $CH = 3.2$。验算：
$BH + HC = 1.8 + 3.2 = 5 = BC$ --- 这是必然的，因为 $H$
位于 $B$ 与 $C$ 之间的斜边上。

\textbf{5.} 两边各除以 $BC$：
\[
BH = \frac{AB^2}{BC},
\qquad
CH = \frac{AC^2}{BC} .
\]
由此可知 $BH$ 与 $CH$ 分别与 $AB^2$ 和 $AC^2$ 成正比例：高的垂足
把斜边分成两段，两段之比等于两直角边的平方之比
（第 4 题中为 $1.8 : 3.2 = 9 : 16$）。

\textbf{6.} 将两式相加，并提取公因式 $BC$
（用到分配律）：
\[
AB^2 + AC^2 = BH \times BC + CH \times BC
= (BH + CH) \times BC = BC \times BC = BC^2 .
\]
这正是 \cref{thm:g8:pythagoras:direct}，仅由余弦关系得出 ---
它是一种真正不同的证明，有别于
\cref{exo:g8:pythagoras:11} 中的四个三角形面积拼图。

\textbf{7.} 在 $ABH$ 中应用毕达哥拉斯定理（以 $H$ 为直角，斜边为
$[AB]$）：$AH^2 = AB^2 - BH^2$。将 $AB^2$ 替换为
$BH \times BC$（第 2 题），并提取公因式 $BH$：
\[
AH^2 = BH \times BC - BH^2
= BH \times (BC - BH)
= BH \times HC .
\]

\textbf{8.} 以两直角边分别为底与高，$ABC$ 的面积
为 $\frac{AB \times AC}{2}$；以斜边 $[BC]$ 为底时，高恰为
$AH$，故面积同样为
$\frac{BC \times AH}{2}$（\cref{thm:g7:areas:triangle}）。两式相等，
两边同乘 2，得 $AB \times AC = BC \times AH$。

\textbf{9.} 故 $AH = \frac{AB \times AC}{BC} = \frac{3 \times 4}{5}
= 2.4$。第 7 题的验证：
$AH^2 = 2.4^2 = 5.76$，且
$BH \times HC = 1.8 \times 3.2 = 5.76$。两者相等。

\textbf{10.} $AH^2 = BH \times HC = 2 \times 8 = 16$，故
$AH = \sqrt{16} = 4$ 厘米。

\textbf{11.} 点 $A$ 位于以 $[BC]$ 为直径的圆上
（且异于 $B$、$C$），由
\cref{thm:g8:cosine:circle} 的第 2 条，三角形 $ABC$ 在
$A$ 处为直角。由作图可知 $(AH) \perp (BC)$，且 $H$ 位于 $B$ 与
$C$ 之间：线段 $[AH]$ 正是第 7 题中的高，
且 $AH^2 = BH \times HC = p \times q$。

\textbf{12.} $AH^2 = 1 \times 2 = 2$：图中仅用直尺和圆规即可作出
一条平方为 $2$ 的线段 --- 即
单位正方形的对角线 $\sqrt2 \approx 1.414$，见
\cref{ex:g8:pythagoras:diagonal}。

\textbf{13.} 作法：首尾相接地放置两条线段，长度分别为
$p = 1$ 与 $q = n$；以整条线段为直径作半圆
（长为 $n + 1$）；在连接点处作垂线，
它与半圆相交于一点，该点
到连接点的距离为 $\sqrt{1 \times n} = \sqrt n$。若要作
$\sqrt5$，则取 $p = 1$ 与 $q = 5$（直径为
$6$ 的半圆）。

\textbf{14.} 点 $A$ 在圆心为 $O$ 的圆上，故
$OA = \frac{p+q}{2}$，即半径。若 $H \neq O$，则三角形
$AHO$ 以 $H$ 为直角，斜边为 $[OA]$，而直角边
短于斜边：$AH < OA$。若 $H = O$，则
$AH = OA$，恰好相等。总之 $AH \leq \frac{p+q}{2}$，
等号当且仅当 $H$ 为圆心时成立 --- 即
$p = q$。

\textbf{15.} 将 $AH \leq \frac{p+q}{2}$ 两边平方（两边均为
正数），并利用 $AH^2 = pq$：
\[
p \times q \leq \left(\frac{p+q}{2}\right)^2 ,
\]
等号当且仅当 $p = q$ 时成立。代数上的再证明借助
乘法公式（\cref{pb:g8:equations:1}）：
\[
\left(\frac{p+q}{2}\right)^2 - pq
= \frac{p^2 + 2pq + q^2 - 4pq}{4}
= \frac{p^2 - 2pq + q^2}{4}
= \left(\frac{p-q}{2}\right)^2 ,
\]
这是一个平方，因此永不为负 --- 并且仅当
$p = q$ 时为零。两个数的几何平均数永远不超过它们
和的一半。
\end{solution}
